\documentclass[a4paper,12pt]{article}
\usepackage[utf8]{inputenc}
\usepackage[T1]{fontenc}
\usepackage[a4paper, top=2cm, bottom=2cm, left=2.5cm, right=2.5cm]{geometry}
\usepackage{booktabs}
\usepackage{array}
\usepackage{listings}
\usepackage{xcolor}
\usepackage{enumitem}
\usepackage{parskip}
\usepackage{titlesec}

\setlength{\parindent}{0pt}
\setlength{\parskip}{4pt}

\titleformat{\section}{\large\bfseries}{}{0em}{}[\vspace{-4pt}\rule{\linewidth}{0.5pt}\vspace{2pt}]
\titleformat{\subsection}{\normalsize\bfseries}{}{0em}{}
\titlespacing{\section}{0pt}{10pt}{4pt}
\titlespacing{\subsection}{0pt}{6pt}{2pt}

\lstset{
  basicstyle=\ttfamily\small,
  backgroundcolor=\color{gray!10},
  frame=single,
  breaklines=true,
  columns=fullflexible,
  keepspaces=true
}

\begin{document}

\begin{center}
  {\LARGE\bfseries Projekt skryptu bashowego}\\[6pt]
  {\large Laboratorium --- Systemy Operacyjne}\\[12pt]
  \begin{tabular}{ll}
    \textbf{Autor:}         & Filip Zdrojewski \\
    \textbf{Nazwa skryptu:} & \texttt{backuper.sh} \\
    \textbf{Data:}          & maj 2026 \\
  \end{tabular}
\end{center}

\vspace{6pt}\hrule\vspace{8pt}

% ============================================================
\section{Opis funkcjonalnosci}

Skrypt \texttt{backuper.sh} realizuje automatyczne, szyfrowane kopie zapasowe
wybranych katalogow uzytkownika z mozliwoscia wysylki na zdalny serwer przez SSH.
Skrypt umozliwia tworzenie kopii, przegladanie listy archiwow oraz przywracanie
danych z wybranej wersji.

\subsection*{Glowne funkcje:}
\begin{itemize}[leftmargin=1.5em, itemsep=1pt, topsep=2pt]
  \item Tworzenie skompresowanych archiwow \texttt{tar.gz} wskazanych katalogow.
  \item Szyfrowanie archiwum kluczem symetrycznym AES-256 przy uzyciu \texttt{gpg}.
  \item Wysylka zaszyfrowanego archiwum na serwer zewnetrzny przez \texttt{scp}/\texttt{rsync}.
  \item Automatyczna rotacja kopii --- przechowywanie ostatnich \textit{N} wersji (konfigurowalne).
  \item Zapis logu operacji (data, katalog, rozmiar, status) do pliku dziennika.
  \item Przywracanie danych z wybranej kopii zapasowej.
  \item Interaktywny wybor kopii do przywrocenia przy uzyciu \texttt{dialog}.
\end{itemize}

% ============================================================
\section{Opis dzialania}

\subsection{Interfejs uzytkownika}

Skrypt obslugiwany jest z linii polecen z opcjami przetwarzanymi przez \texttt{getopts}.
Operacje wymagajace wyboru (przywracanie) korzystaja z interfejsu \texttt{dialog}.

\begin{lstlisting}
Uzycie: backuper.sh [OPCJE]

  -b            wykonaj kopie zapasowa (backup)
  -r            przywroc dane z wybranej kopii (tryb dialog)
  -l            wyswietl liste dostepnych kopii na serwerze
  -d KATALOG    katalog zrodlowy do backupu (mozna wielokrotnie)
  -o KATALOG    katalog docelowy przywracania (domyslnie: ~/przywrocone)
  -n N          liczba przechowywanych kopii (domyslnie z pliku .rc)
  -v            wersja i autor
  -h            krotka pomoc
\end{lstlisting}

\subsection{Plik konfiguracyjny \texttt{backuper.rc}}

Dane konfiguracyjne przechowywane sa w pliku \texttt{\char126/.backuper.rc}:

\begin{lstlisting}
BACKUP_DIRS="/home/user/dokumenty /home/user/projekty"
REMOTE_HOST="backup@192.168.1.100"
REMOTE_DIR="/backups/filip"
KEEP_LAST=7
GPG_PASS_FILE="~/.backup_pass"
LOG_FILE="~/.backuper.log"
\end{lstlisting}

\subsection{Przebieg operacji backup}

\begin{enumerate}[leftmargin=1.5em, itemsep=1pt, topsep=2pt]
  \item Odczyt konfiguracji z pliku \texttt{.rc}.
  \item Utworzenie archiwum \texttt{tar.gz} w \texttt{/tmp} (dane tymczasowe usuwane przez \texttt{trap}).
  \item Zaszyfrowanie: \texttt{gpg --symmetric --cipher-algo AES256}.
  \item Wysylka pliku \texttt{.tar.gz.gpg} na serwer przez \texttt{scp}.
  \item Rotacja --- usuniecie nadmiarowych kopii przez \texttt{ssh} + \texttt{find}.
  \item Zapis wpisu do logu i czyszczenie \texttt{/tmp}.
\end{enumerate}

\subsection{Przebieg operacji przywracania}

\begin{enumerate}[leftmargin=1.5em, itemsep=1pt, topsep=2pt]
  \item Pobranie listy kopii z serwera (\texttt{ssh} + \texttt{ls}).
  \item Wyswietlenie listy w interfejsie \texttt{dialog --menu}.
  \item Pobranie wybranego archiwum przez \texttt{scp} do \texttt{/tmp}.
  \item Odszyfrowanie: \texttt{gpg --decrypt}.
  \item Rozpakowanie do katalogu docelowego (\texttt{tar -xzf}).
  \item Usuniecie plikow tymczasowych z \texttt{/tmp}.
\end{enumerate}

% ============================================================
\section{Narzedzia i technologie}

\begin{center}
\begin{tabular}{>{\ttfamily}lp{9cm}}
  \toprule
  \normalfont\textbf{Narzedzie} & \textbf{Zastosowanie} \\
  \midrule
  tar       & Tworzenie i rozpakowywanie archiwow skompresowanych \\
  gpg       & Szyfrowanie i deszyfrowanie symetryczne AES-256 \\
  scp       & Transfer plikow na zdalny serwer przez SSH \\
  ssh       & Zdalne zarzadzanie kopiami (rotacja, listowanie) \\
  find      & Wyszukiwanie i usuwanie nadmiarowych kopii \\
  dialog    & Interaktywny wybor kopii do przywrocenia \\
  getopts   & Obsluga opcji linii polecen \\
  trap      & Czyszczenie plikow tymczasowych przy wyjsciu \\
  \bottomrule
\end{tabular}
\end{center}

% ============================================================
\section{Wymagania}

\begin{itemize}[leftmargin=1.5em, itemsep=1pt, topsep=2pt]
  \item System operacyjny: Linux (bash $\geq$ 4.0)
  \item Pakiety: \texttt{gpg}, \texttt{dialog}, \texttt{openssh-client}, \texttt{tar}
  \item Dostep SSH do serwera zewnetrznego (klucz lub haslo)
  \item Skonfigurowany plik \texttt{\char126/.backuper.rc}
\end{itemize}

\end{document}
